% ID: SEL-D0035-OB-E04
% ESTADO: cotejada. Segunda pasada visual y borrador compilado revisados.
% Fuente: Selectividad/matematicas_aplicadas/titular_septiembre_EXAMEN 1-ANDALUCÍA.pdf
% SHA-256: 7a09a02641ea8de345b94f49b27515e58436e032c7401867f86dc523490a3ea5
% Página: 2; opción B; ejercicio original 4.

En una muestra, elegida al azar, de 100 estudiantes de una Universidad, se ha observado que 25 desayunan en la cafetería del campus.
\begin{enumerate}[label=\alph*)]
\item \textbf{(1.25 puntos)} Determine, con un nivel de confianza del 95\%, un intervalo de confianza para estimar la proporción de estudiantes de esa Universidad que desayunan en la cafetería.
\item \textbf{(1.25 puntos)} Si la proporción de estudiantes de esa Universidad que desayunan en la cafetería del campus en una muestra aleatoria es de 0.2, y el error cometido en la estimación ha sido inferior a 0.03, con un nivel de confianza del 92.5\% calcule el tamaño mínimo de la muestra.
\end{enumerate}
